\section{Human-preference extension: HelpSteer2 methods}
\label{sec:helpsteer}
\noindent\textbf{Protocol scope.} This section preserves the initial
HelpSteer2 reward-model and unigram-RL design. The later semantic-MMD
profile in Section~\ref{sec:mmd} uses a different frozen scorer,
reference subset, generation limit, and conditional batching. Its
settings must not be inferred from this earlier design.

\noindent\textbf{Status of evidence.} This section records the original
prepared-data and reward-model design, including the larger configured
policy profile. Section~\ref{sec:newresults} reports the completed subset
trials and subsequent corrected profiles with their actual budgets and
scorer-selection audit. Preparation counts below are not achieved
policy-training counts; later configurations are stated separately.

\subsection{Dataset and prompt-disjoint preparation}
We replace a synthetic mathematical target with natural-language prompts
and human feedback from \texttt{nvidia/HelpSteer2}
\cite{wang2024helpsteer,wang2024preference}. We use both the attribute-rated
responses and the separate \texttt{preference/preference.jsonl.gz} pairs,
pinned to dataset revision \texttt{990b2711a361}. The full revision and
source-file hashes appear in the supplementary protocol record. The
dataset is distributed under CC-BY-4.0; its preference annotations are
human feedback, not synthetic judgments generated for this experiment.

Before splitting, the preparation code collapses whitespace in the prompt
text and hashes the result. Whitespace-equivalent prompts therefore stay
together. One pair is retained per canonical prompt, and an official
validation occurrence takes priority over a duplicate training occurrence.
This grouping does not establish semantic deduplication of differently
worded prompts. Full multi-turn context is parsed from the source
User/Assistant markers and preserved in the Qwen chat template.

The original training partition is divided by a deterministic prompt-hash
rule: hash modulo four equal to zero reserves a prompt for the policy/PMF
pool, and the other residues reserve it for reward-model training.
The nominal allocation is 25\%/75\% before further policy eligibility
filters, not an exact fraction after filtering. The original validation
partition supplies disjoint reward validation, policy validation, and
policy test prompts. Eligible policy evaluation prompts are sorted,
shuffled with seed 42, and assigned to validation and test; the remaining
original validation prompts supply reward validation.

Both candidate responses must have at least 90\% ASCII among alphabetic
characters. This is a language-selection heuristic rather than an English
classifier. Chat-formatted prompts must have at most 384 tokens, and
each complete candidate reward-model input at most 1,024 tokens. Source
answers exceeding a limit are excluded rather than silently truncated.
Policy/reference eligibility additionally requires a non-tied preferred
answer, at most 256 response tokens, with human coherence, correctness,
and helpfulness ratings each at least 3 on the 0--4 scale. Reward-model
training does not require this additional preferred-reference filter.
These choices define a selected short-answer subset, not the complete
HelpSteer2 benchmark.

\begin{table}[htbp]
\centering\small
\caption{Verified prepared HelpSteer2 subset. Splits are disjoint by canonical
prompt hash. These counts do not indicate completed training or evaluation.}
\label{tab:helpsteerdata}
\begin{tabular}{@{}lrp{0.52\linewidth}@{}}
\toprule
Subset & Count & Use\\
\midrule
Reward training & 5,095 pairs & Fit the scalar preference reward model\\
Reward validation & 295 pairs & Select reward checkpoint and normalize scores\\
Policy training & 533 prompts & Sample actions for both policy arms\\
Preferred-response reference & 533 responses & Build the training-only word PMF\\
Policy validation & 16 prompts & Periodic policy diagnostics\\
Policy test & 32 prompts & Final base/quality/KL comparison\\
\bottomrule
\end{tabular}
\end{table}

The reward-training and reward-validation sets include 1,225 and 79 ties,
respectively; reward validation therefore has 216 non-tied pairs.
The reference PMF uses only the 533 preferred policy-training responses.
Reward-training and policy-training prompts are disjoint; evaluation
responses and labels enter neither reward training nor PMF construction.
Held-out human ratings do affect evaluation-subset eligibility, which must
be considered when interpreting eventual results. Human-written preference
rationales are not inputs to the trained reward scorer. Neither the prior
Shakespeare corpus nor a Shakespeare-trained adapter supplies data or
initialization for this extension.

\subsection{Bradley--Terry preference reward model}
Let a human-labeled pair be $(c,y_1,y_2,s)$, where positive source preference
strength $s$ prefers \emph{response two}, negative $s$ prefers response one,
and zero is a tie. A separate
\texttt{Qwen2.5-3B-Instruct} sequence classifier produces a scalar score
$r_\phi(c,y)$. With $\sigma(z)=(1+e^{-z})^{-1}$ and
$\Delta_\phi=r_\phi(c,y_2)-r_\phi(c,y_1)$, the pairwise model is
\begin{equation}
 P_\phi(y_2\succ y_1\mid c)=\sigma(\Delta_\phi),
 \qquad t(s)=\frac{1+\operatorname{sign}(s)}{2}.
 \label{eq:btprob}
\end{equation}
We minimize binary cross-entropy with this soft target:
\begin{align}
 \mathcal L_{\mathrm{BT}}(\phi)
 &=-\E_{\mathcal D_{\mathrm{RM}}}
  [t\log\sigma(\Delta_\phi)+(1-t)\log(1-\sigma(\Delta_\phi))],
 \label{eq:btloss}\\
 \frac{\partial\mathcal L_{\mathrm{pair}}}{\partial\Delta_\phi}
 &=\sigma(\Delta_\phi)-t.
 \label{eq:btgradient}
\end{align}
The sign convention matters: a positive label should increase the second
answer's score relative to the first. Preference-strength magnitude is
not interpreted as a probability or a loss weight. Ties use $t=1/2$,
whose pair loss is minimized at equal scores. This is soft-label treatment
of ties, not a separate three-outcome tie-probability model. The pair
likelihood depends on score differences and is invariant to a shared
additive score shift.

The reward model has attention LoRA rank 16, alpha 32, no dropout, and a
trainable scalar score head. The planned fit is one shuffled epoch over
5,095 pairs, effective batch 32 pairs, microbatch four pairs, AdamW learning
rate $2\times10^{-5}$, cosine decay, approximately 5\% warmup, weight
decay 0.01, and gradient-norm clipping at one. A full epoch would contain
160 updates; the elapsed-time budget can stop it earlier. Model selection
uses reward-validation pair loss. Before policy training, the selected
model must have non-tied validation accuracy greater than 50\%, validation
score standard deviation above $10^{-6}$, and at least one training
update. This is a technical gate, not a significance test or certification
of generated-answer quality. A subsequent stronger selection audit and
its outcomes appear in Section~\ref{sec:newresults}.

For the selected reward model, pool both candidates' raw scores on reward
validation and fix their mean $m_\phi$ and population standard deviation
$s_\phi$. The frozen quality reward for a policy action is
\begin{equation}
 r_Q(c,y)=\ind\{y\text{ terminates with EOS}\}\,
 \sigma\!\left(\frac{r_\phi(c,y)-m_\phi}{\max(s_\phi,10^{-6})}\right).
 \label{eq:qualitytransform}
\end{equation}
The implementation clips the standardized score to $[-30,30]$ before
applying the sigmoid. This bounded utility is not a calibrated human
win probability; centering and scaling do not supply that guarantee.
The scorer sees the complete prompt and answer, without human rationales
at inference. Its parameters and normalization remain fixed during RL.

\subsection{Training-only word target and global histogram}
Let $\mathcal R$ be the 533 preferred policy-training responses. Words
are extracted using ASCII alphabetic runs with optional internal straight
or curly apostrophes, lowercased, and curly apostrophes normalized to
straight ones. No stemming or stopword removal is applied. The vocabulary
$\mathcal V$ consists only of reference words. A single pooled unigram
target is
\begin{equation}
 p_w=\frac{C_{\mathcal R}(w)}{\sum_{v\in\mathcal V}C_{\mathcal R}(v)},
 \quad w\in\mathcal V,\qquad p_{\mathrm{UNK}}=0.
 \label{eq:wordtarget}
\end{equation}
The prepared reference contains 57,251 extracted word occurrences and
7,593 distinct word types. These are corpus statistics, not policy
measurements. This is a marginal word-frequency target across responses, not a
prompt-conditional law or a measure of semantic correctness.

For a global batch $\mathcal B=(Y_1,\ldots,Y_M)$, collect each passage's
word-count vector $\mathbf c_i$ on $\mathcal V\cup\{\mathrm{UNK}\}$.
Unseen words enter UNK. An empty extracted passage contributes one UNK
count; unfinished passages retain all their observed words. Let
$V=|\mathcal V|$, $C_w=\sum_i c_{i,w}$, and $L=\sum_w C_w$. The
total smoothing mass is fixed at $a=1$, unlike the numerical experiment's
per-category pseudocount 0.1:
\begin{equation}
 \widehat q_w=\frac{C_w+a/(V+1)}{L+a},\qquad
 D_{\mathrm{word}}(\mathcal B)=
 \sum_{w:p_w>0}p_w\log\frac{p_w}{\widehat q_w}.
 \label{eq:wordkl}
\end{equation}
The denominator includes UNK counts. Because counts pool words, long
passages contribute more mass, and even a perfectly matched population
PMF has sample-count-dependent estimated KL. Future comparisons must
report lengths and use count-matched post-hoc diagnostics; a lower finite
KL alone is not sufficient evidence of better lexical calibration.

\subsection{Controlled policy arms and the combined objective}
The two independent actor adapters start from the same pinned
\texttt{Qwen2.5-3B-Instruct} checkpoint and seed 42. One arm uses quality
reward only, $\lambda=0$; the other adds word-distribution control,
$\lambda=0.1$. Both use the same frozen learned reward and training
prompts. The unadapted instruction checkpoint, evaluated with actor
adapters disabled, is the third comparison condition. Unlike numerical
RL, this stage does not initialize from the numerical SFT adapter.

Writing $\pi_0$ for the unadapted instruction policy, the objective is
\begin{align}
 J_\lambda(\theta)=
 &\E_{\mathbf c,Y_{1:M}}
 \left[\frac1M\sum_{i=1}^{M}r_Q(c_i,Y_i)
                    -\lambda D_{\mathrm{word}}(Y_{1:M})\right]\nonumber\\
 &-\beta\E_c\KL\!\left(\pi_\theta(\cdot\mid c)
                         \Vert\pi_0(\cdot\mid c)\right),
 \qquad \beta=0.002.
 \label{eq:helpsteerobjective}
\end{align}
Here $\mathbf c$ is the batch of prompts. The Bradley--Terry likelihood
fits the reward model offline; on-policy REINFORCE updates the actor
later. These are different training stages. The numeric target KL,
word-target KL, and reference-policy KL also remain distinct quantities.

For $D=D_{\mathrm{word}}(\mathcal B)$ and
$D_{-i}=D_{\mathrm{word}}(\mathcal B\setminus\{Y_i\})$, the global
leave-one-passage-out credits are
\begin{align}
 A_i^Q&=r_Q(c_i,Y_i)-\frac{1}{M-1}\sum_{j\ne i}r_Q(c_j,Y_j),\\
 A_i^D&=M(D_{-i}-D),\\
 \widehat{\mathcal L}_{\mathrm{actor}}&=
 -\frac1M\sum_i\sg\!\left[A_i^Q+\lambda A_i^D
       -\beta\log\frac{\pi_\theta(Y_i\mid c_i)}{\pi_0(Y_i\mid c_i)}\right]
       \log\pi_\theta(Y_i\mid c_i).
 \label{eq:helpsteersurrogate}
\end{align}
Removing a passage removes its entire count vector and changes the
histogram's word-count denominator. Advantages are computed on the
global batch gathered from both GPU ranks, not on separate per-rank
histograms. Equal local batches and DDP averaging yield the global mean
gradient. Quality and weighted distribution-advantage RMS values are
logged separately; no within-batch variance normalization rescales them.

\begin{table}[htbp]
\centering\small
\caption{Original larger HelpSteer2 configuration, not the achieved subset
trial. Completed-run counts and overrides appear in Section~\ref{sec:newresults}.}
\label{tab:helpsteersettings}
\begin{tabular}{@{}lp{0.62\linewidth}@{}}
\toprule
Setting & Configuration\\
\midrule
Actor budget & Up to 600 updates, 16 global passages/update (9,600 if complete)\\
Prompt grouping & Two prompts/update, eight completions per prompt across two GPUs\\
Actor adaptation & Attention LoRA rank 16, alpha 32, dropout disabled\\
Actor optimizer & AdamW, learning rate $5\times10^{-6}$, weight decay 0.01\\
Schedule & Cosine decay, 40 warmup updates, gradient clipping at one\\
Gradient microbatch & Two passages per GPU microbatch\\
Sampling & Temperature 1, top-$p=1$, top-$k$ disabled, at most 384 new tokens\\
Word-KL weight & $\lambda\in\{0,0.1\}$, not tuned for the new learned reward\\
Policy validation / test & 16 / 32 prompts, two draws per prompt\\
Reference penalty & Reverse sequence KL to instruction base, coefficient 0.002\\
\bottomrule
\end{tabular}
\end{table}

\subsection{Original evaluation design and interpretation}
The configured policy validation uses 32 passages per probe and final
test uses 64 passages per condition. Initial and final policies use
matched per-rank seeds and decoding settings. A separate pinned
Qwen2.5-7B-Instruct auditor rates coherence, readability, and prompt
fulfillment on final test generations. The auditor never supplies the
training reward. It computes expected ratings conditional on tokens
1--5, records their total original probability, and rejects a rating
whose token mass is below 0.01. It is a separate same-family model,
not human evaluation or an independently validated quality standard.

The intended analysis compares the instruction baseline, quality-only
actor, and quality-plus-KL actor on learned reward, external rubric
scores, word KL, UNK rate, EOS termination, and length. Actual training
counts and stop reasons must be included: an elapsed-time cutoff can
leave either arm below 600 updates, weakening an equal-compute
interpretation. The reward fit has a 55-minute elapsed cutoff within a
one-H200 one-hour allocation; each sequential policy job has a
105-minute training cutoff within a two-H200 two-hour allocation,
reserving time for evaluation. These are configured limits, not measured
runtimes. The recorded subset trials and later corrected runs in
Section~\ref{sec:newresults} supersede this profile's pending status;
the full 600-update design is not a completed result.

Training a reward from human labels makes its supervision human-derived;
it does not make generated policy outputs human-validated. A future
positive result would require improvements that survive length/count
controls and do not merely exploit the learned scorer or its word
reference. Multi-seed comparisons and human evaluation would still be
needed. The completed trials do not establish a reliable joint gain.

The source archive's \texttt{supplementary/helpsteer-protocol.json} records
the full dataset and model revisions, prepared-file checksums, split and
tie counts, reference-vocabulary statistics, config, and implementation
hashes. Completed performance records are supplied separately under
\texttt{supplementary/new-results/}.
